\chapter{Hip Replacement}

\section*{What Is This Surgery?}
Total hip arthroplasty (THA), commonly referred to as hip replacement surgery, is a significant orthopedic intervention aimed at excising a diseased or damaged hip joint and reconstructing it with prosthetic components. This procedure is primarily indicated for patients suffering from debilitating conditions such as severe osteoarthritis, rheumatoid arthritis, or avascular necrosis, which lead to chronic pain and functional impairment. The surgery involves the replacement of the femoral head with a metallic femoral stem, which is inserted into the femur, and an acetabular cup that is secured into the pelvic socket. The primary goal of THA is to alleviate pain, restore mobility, and enhance the overall quality of life for patients who have not responded to conservative treatments. This surgery is characterized by its high success rates and significant improvements in patient-reported outcomes.

\section*{Alternative Names}
\begin{itemize}
  \item Total hip arthroplasty (THA)
  \item Total hip replacement surgery
  \item Hip joint replacement
  \item Prosthetic hip replacement
  \item Hip reconstruction surgery
  \item Hip resurfacing (in specific cases)
\end{itemize}

\section*{Relevant Surgical Anatomy}
The hip joint is a ball-and-socket synovial joint formed by the articulation of the head of the femur and the acetabulum of the pelvis. The acetabulum is a cup-shaped depression located on the lateral aspect of the pelvis, formed by the fusion of the ilium, ischium, and pubis. The femoral head is covered by articular cartilage, which, along with the acetabulum, facilitates smooth, low-friction movement. 

Key ligaments that provide stability to the hip joint include the iliofemoral, pubofemoral, and ischiofemoral ligaments. The ligamentum teres, which contains a small artery supplying the femoral head, is also present within the joint. 

Surrounding muscles play a crucial role in hip function and are often encountered during surgical approaches. These include:
\begin{itemize}
  \item Gluteal muscles (maximus, medius, minimus) for abduction and extension.
  \item Iliopsoas for hip flexion.
  \item Adductor group for adduction.
\end{itemize}

Major neurovascular structures in proximity to the hip joint include the femoral nerve, artery, and vein anteriorly, the sciatic nerve posteriorly, and the obturator nerve and vessels medially. The profunda femoris artery and its circumflex branches provide the primary blood supply to the femoral head and neck, while venous drainage parallels the arterial supply. Lymphatic drainage primarily flows to the superficial and deep inguinal lymph nodes. Surgical landmarks such as the anterior superior iliac spine, greater trochanter, and pubic symphysis are critical for guiding incision placement and component alignment.

\section*{Surgical Principle \& Rationale}
The fundamental principle of total hip replacement is to excise the diseased or damaged articular surfaces of both the femoral head and the acetabulum, replacing them with durable, biocompatible prosthetic components. In conditions like severe osteoarthritis, the protective articular cartilage erodes, leading to painful bone-on-bone friction, inflammation, and loss of joint mobility. The rationale is to eliminate this painful articulation by creating a new, smooth, low-friction bearing surface using metal, ceramic, and highly cross-linked polyethylene materials.

The technical goals of the procedure include restoring the anatomical center of rotation of the hip, correcting any existing leg length discrepancy, achieving a stable joint with an optimal range of motion to prevent dislocation, and ensuring robust, long-lasting fixation of the prosthetic components to the host bone. Fixation can be achieved either through bone cement (cemented fixation), which provides immediate stability, or by a press-fit technique (cementless fixation), which relies on biological ingrowth of bone into a porous implant surface over time. This comprehensive reconstruction aims to provide significant pain relief and functional restoration, allowing patients to resume daily activities with improved comfort and mobility.

\section*{History \& Surgical Evolution}
The concept of hip arthroplasty dates back to the 19th century, with early attempts involving interpositional arthroplasty using various materials. However, modern total hip replacement was revolutionized by Sir John Charnley, a British orthopedic surgeon, in the early 1960s. Charnley's pioneering work at Wrightington Hospital established the principles of low-friction arthroplasty, utilizing a small-diameter metallic femoral head articulating with a high-density polyethylene acetabular cup, both fixed to the bone with acrylic bone cement. His meticulous surgical technique, emphasis on sterile operating environments, and standardized implant design laid the foundation for the procedure's widespread success and longevity.

Following Charnley's innovations, the field continued to evolve. In the 1970s and 1980s, cementless implants were developed, featuring porous surfaces designed to encourage bone ingrowth, offering an alternative fixation method, particularly for younger, more active patients. The 1990s saw advancements in bearing surfaces, including highly cross-linked polyethylene, which significantly reduced wear rates, and ceramic-on-ceramic or ceramic-on-polyethylene options, further improving implant durability. More recently, minimally invasive surgical approaches and the anterior approach have gained popularity, aiming to reduce soft tissue dissection, minimize postoperative pain, and accelerate recovery, though their long-term benefits are still debated.

\section*{Clinical Indications}
Total hip replacement is indicated in several clinical scenarios, including:

\begin{itemize}
  \item Severe osteoarthritis of the hip, characterized by debilitating pain, stiffness, and functional limitation unresponsive to conservative management.
  \item Rheumatoid arthritis or other inflammatory arthropathies causing destructive changes to the hip joint.
  \item Avascular necrosis (osteonecrosis) of the femoral head, leading to collapse of the femoral head and severe pain.
  \item Post-traumatic arthritis resulting from previous hip fractures or dislocations, causing secondary degenerative changes.
  \item Certain displaced femoral neck fractures in elderly patients, where internal fixation is unlikely to succeed or is contraindicated.
  \item Developmental dysplasia of the hip (DDH) in adulthood, leading to premature degenerative changes.
  \item Benign or malignant tumors of the proximal femur or acetabulum requiring resection and reconstruction.
  \item Failed previous hip surgeries, such as osteotomies or hemiarthroplasties, resulting in persistent pain and dysfunction.
\end{itemize}

\section*{Clinical Positioning}
Total hip replacement is predominantly positioned as a primary, definitive surgical treatment for end-stage hip joint disease. It is considered the gold standard intervention when non-surgical modalities, including physical therapy, anti-inflammatory medications, activity modification, and corticosteroid injections, have failed to provide adequate pain relief or functional improvement. In this context, it serves as a curative procedure aimed at restoring joint integrity and function.

While primarily a first-line definitive treatment for severe arthritis, THA can also function as a secondary or salvage procedure. This applies to situations where previous surgical interventions, such as internal fixation of a hip fracture, hip osteotomy, or even a prior hemiarthroplasty, have failed due to non-union, avascular necrosis, or progressive arthritis. In these cases, THA is performed to address the complications of the initial surgery and provide a durable solution. The decision to proceed with THA is always made after careful consideration of the patient's symptoms, functional limitations, radiographic findings, and overall health status.

\section*{Surgical Technique \& Procedure}
Total hip replacement is typically performed under general anesthesia or spinal/epidural anesthesia with sedation. The patient is positioned either supine (for anterior approach) or in a lateral decubitus position (for posterior or lateral approaches), with careful padding to prevent nerve compression. Prophylactic broad-spectrum antibiotics are administered intravenously prior to incision, and the surgical site is meticulously prepared and draped in a sterile fashion.

The key operative steps generally include:

\begin{itemize}
  \item **Incision and Approach:** A skin incision is made, and the surgeon proceeds through muscle layers to expose the hip joint. Common approaches include the posterior (preserving abductors, risk of dislocation), lateral (splitting abductors, less dislocation risk), and anterior (muscle-sparing, often minimally invasive).
  
  \item **Femoral Head Resection:** The hip joint is dislocated, and the femoral neck is osteotomized (cut) at a predetermined level, allowing removal of the diseased femoral head.
  
  \item **Acetabular Preparation:** Specialized reamers are used to remove the remaining articular cartilage and subchondral bone from the acetabulum, creating a hemispherical cavity to accept the prosthetic cup.
  
  \item **Acetabular Component Implantation:** The acetabular cup, often porous-coated for bone ingrowth, is impacted into the prepared socket. Screws may be used for additional fixation, and a polyethylene or ceramic liner is then inserted into the cup.
  
  \item **Femoral Canal Preparation:** The femoral canal is progressively broached (reamed) to create a precise fit for the femoral stem. The broach size is selected to match the chosen femoral stem.
  
  \item **Femoral Stem and Head Implantation:** The femoral stem is inserted into the prepared femoral canal, either cemented or press-fit. A trial femoral head is placed on the stem, and the joint is reduced to assess stability, leg length, and range of motion. Once optimal parameters are achieved, the definitive femoral head (metal or ceramic) is affixed.
  
  \item **Joint Reduction and Stability Check:** The new prosthetic joint is carefully reduced, and the surgeon performs a series of maneuvers to confirm stability, assess leg length equality, and ensure an adequate range of motion without impingement or dislocation.
  
  \item **Wound Closure:** The surgical site is irrigated, and the deep tissues, muscles, and fascia are meticulously repaired in layers. The subcutaneous tissue and skin are closed with sutures or staples. A surgical drain may be placed temporarily to collect any postoperative bleeding.
\end{itemize}

The entire procedure typically takes 1 to 2 hours, after which the patient is transferred to the post-anesthesia care unit (PACU) for initial recovery.

\section*{Preoperative Workup \& Preparation}
Thorough preoperative evaluation is critical to optimize patient outcomes and minimize risks. This includes a comprehensive medical history, physical examination, and assessment of comorbidities. Routine laboratory tests typically include a complete blood count, basic metabolic panel, coagulation profile, and urinalysis. An electrocardiogram (ECG) and chest X-ray are standard, especially for older patients or those with cardiac or pulmonary history. Cardiac and anesthesia clearance are obtained, often involving consultation with an internist or cardiologist, to assess surgical risk (e.g., using the ASA physical status classification system).

Imaging studies include standard anteroposterior and lateral radiographs of the hip and pelvis, which are used for templating to plan component sizes, alignment, and identify any anatomical variations. Patients are advised to stop certain medications, particularly anticoagulants and antiplatelet agents, several days to a week prior to surgery, as directed by the surgeon and anesthesiologist. Smoking cessation is strongly encouraged to improve wound healing and reduce complication rates. Dental work should be completed before surgery to mitigate the risk of hematogenous spread of infection to the prosthetic joint. Preoperative physical therapy may be recommended to strengthen surrounding muscles and educate the patient on postoperative exercises and precautions. On the day of surgery, patients are instructed to fast for at least 6-8 hours, and prophylactic antibiotics are administered. Informed consent, detailing the procedure, risks, benefits, and alternatives, is obtained.

\section*{Contraindications \& Precautions}
Absolute contraindications to total hip replacement include:

\begin{itemize}
  \item Active systemic infection (e.g., sepsis) or active local infection in the hip joint or surrounding tissues.
  \item Severe neurological deficits that render the limb non-functional or prevent safe ambulation and rehabilitation.
  \item Rapidly progressive neurological disease that would preclude effective rehabilitation.
  \item Skeletal immaturity (open growth plates), except in rare, specific circumstances.
  \item Charcot arthropathy of the hip, due to high risk of implant loosening and failure.
  \item Uncontrolled medical comorbidities that pose an unacceptably high anesthetic or surgical risk.
\end{itemize}

Relative contraindications and precautions include:

\begin{itemize}
  \item Poorly controlled chronic medical conditions (e.g., diabetes, hypertension, chronic kidney disease) that significantly increase surgical and anesthetic risks.
  \item Morbid obesity, which can increase surgical complexity, complication rates (infection, dislocation), and implant wear.
  \item Active infection elsewhere in the body (e.g., urinary tract infection, dental abscess) that should be treated prior to surgery.
  \item Severe osteoporosis or extremely poor bone quality, which may compromise implant fixation.
  \item Significant peripheral vascular disease or severe venous insufficiency, increasing the risk of wound complications and deep vein thrombosis.
  \item Unrealistic patient expectations regarding the surgical outcome, activity levels, or pain relief.
  \item Non-compliance with preoperative instructions or anticipated poor adherence to postoperative rehabilitation protocols.
\end{itemize}

\section*{Complications \& Risks}
While total hip replacement is generally safe and highly effective, like any major surgery, it carries potential complications and risks. These can be broadly categorized into general surgical risks and procedure-specific risks.

General surgical complications include:

\begin{itemize}
  \item **Anesthesia-related risks:** Adverse reactions to anesthetic agents, respiratory depression, cardiac events (myocardial infarction, arrhythmia), stroke.
  \item **Bleeding:** Intraoperative blood loss requiring transfusion, or postoperative hematoma formation.
  \item **Infection:** Superficial wound infection or deep periprosthetic joint infection (PJI), which can be devastating and often requires further surgery.
  \item **Thromboembolic events:** Deep vein thrombosis (DVT) in the leg, which can lead to pulmonary embolism (PE), a potentially life-threatening complication.
\end{itemize}

Procedure-specific risks include:

\begin{itemize}
  \item **Hip dislocation:** The prosthetic femoral head can dislodge from the acetabular cup, particularly in the early postoperative period or with certain movements.
  \item **Leg length discrepancy:** The operated leg may be perceived as longer or shorter than the contralateral limb, which can sometimes require shoe lifts or further intervention.
  \item **Nerve injury:** Damage to nerves such as the sciatic, femoral, or obturator nerves, leading to weakness, numbness, or paralysis in the affected limb.
  \item **Vascular injury:** Damage to major blood vessels, though rare, can lead to significant bleeding or limb ischemia.
  \item **Periprosthetic fracture:** Fracture of the femur or pelvis can occur during implant insertion or postoperatively due to trauma or stress shielding.
  \item **Implant loosening or wear:** Over time, the prosthetic components can loosen from the bone or the bearing surfaces can wear out, necessitating revision surgery.
  \item **Heterotopic ossification:** Abnormal bone formation in the soft tissues around the hip, which can restrict motion.
  \item **Aseptic loosening:** Loosening of the implant without infection, often due to biological reactions to wear debris or mechanical failure.
  \item **Persistent pain:** While rare, some patients may experience ongoing pain despite a technically successful surgery.
\end{itemize}

\section*{Postoperative Recovery Timeline}
Immediately after surgery, the patient is transferred to the Post-Anesthesia Care Unit (PACU) for close monitoring of vital signs, pain levels, and neurological status. Once stable, they are moved to a hospital ward. Early mobilization is a cornerstone of recovery, often beginning on the day of surgery or the following morning, with the assistance of a physical therapist. Patients are taught hip precautions (e.g., avoiding extreme flexion, adduction, and internal rotation) to minimize the risk of dislocation. Pain management is initiated with a multimodal approach, including oral analgesics, nerve blocks, and sometimes patient-controlled analgesia (PCA). Anticoagulation therapy is typically continued for several weeks to prevent deep vein thrombosis.

During the first 24-48 hours, the focus is on pain control, preventing complications, and initiating basic mobility such as sitting up, transferring to a chair, and short walks with a walker or crutches. Hospital stay typically ranges from 1 to 3 days. Over the first 1-2 weeks post-discharge, patients continue with prescribed exercises, gradually increasing walking distances. Diet is advanced as tolerated. Activity restrictions include avoiding heavy lifting, twisting, and high-impact activities. Most patients transition from a walker to crutches or a cane within 2-4 weeks. Long-term recovery involves continued physical therapy for several months to regain full strength, endurance, and range of motion. Return to most normal daily activities, including driving, usually occurs within 4-6 weeks, with full recovery and return to light recreational activities often taking 3-6 months.

\section*{Surgical Findings \& Pathology}
During total hip replacement, the surgeon meticulously documents the intraoperative findings. This typically includes the extent of articular cartilage loss (e.g., eburnation, fibrillation), presence and size of osteophytes (bone spurs), subchondral cysts, and any deformities of the femoral head or acetabulum. The quality of the bone, presence of synovitis, and integrity of the joint capsule and surrounding soft tissues are also noted. Any signs of avascular necrosis, such as collapse or discoloration of the femoral head, are recorded.

The resected femoral head and any excised acetabular osteophytes are sent to pathology for histological examination. The pathology report typically confirms the diagnosis, such as osteoarthritis, rheumatoid arthritis, or avascular necrosis, by describing characteristic microscopic changes in the bone and cartilage. For osteoarthritis, this includes cartilage degeneration, subchondral sclerosis, and osteophyte formation. In cases of avascular necrosis, areas of necrotic bone and marrow are identified. The report also rules out any unexpected findings, such as tumor or infection, which is crucial for definitive diagnosis and guiding any further management.

\section*{Expected Postoperative Course}
A normal and successful postoperative course following total hip replacement is characterized by a progressive reduction in pain, gradual restoration of hip function, and uneventful wound healing. Patients should expect initial incisional pain, which is well-managed with medication, and this pain should steadily decrease over days to weeks. Swelling and bruising around the hip are common and typically resolve within a few weeks. The patient's ability to bear weight and ambulate with assistive devices improves daily, transitioning from a walker to crutches, then a cane, and eventually independent walking.

The hip precautions are diligently followed to prevent dislocation, and adherence to the physical therapy regimen is crucial for regaining strength and mobility. By 6-12 weeks, most patients experience significant pain relief and are able to perform most activities of daily living with minimal assistance. Radiographs taken postoperatively should confirm optimal positioning and fixation of the prosthetic components, with no signs of fracture or dislocation. The ultimate goal is a stable, pain-free hip that allows the patient to return to a functional and active lifestyle.

\section*{Postoperative Care \& Follow-up}
Postoperative care involves a structured rehabilitation program and regular follow-up with the orthopedic surgeon.

\begin{itemize}
  \item **Hospital Stay:** Typically 1-3 days, focusing on pain management, DVT prophylaxis, and early mobilization with physical therapy.
  
  \item **Discharge Planning:** Patients are discharged home or to a rehabilitation facility, with instructions on hip precautions, wound care, medication management, and exercise.
  
  \item **Physical Therapy:** Continues for several weeks to months, focusing on strengthening, range of motion, gait training, and balance.
  
  \item **Follow-up Visits:** Scheduled at approximately 2 weeks (for wound check, staple/suture removal), 6 weeks, 3 months, 6 months, 1 year, and then annually or biennially.
  
  \item **Imaging:** X-rays of the hip are taken at follow-up visits to monitor implant position, alignment, and detect any signs of loosening or wear.
  
  \item **Warning Signs:** Patients are educated to contact their doctor immediately for signs of infection (fever, chills, redness, pus from wound), sudden severe hip pain, inability to bear weight, or signs of DVT (calf pain, swelling, tenderness).
\end{itemize}

Long-term surveillance is essential to monitor implant longevity and address any potential issues promptly.

\section*{Epidemiology}
Total hip replacement is one of the most frequently performed and successful orthopedic procedures globally. The incidence of THA has steadily increased over the past few decades, driven by an aging population, rising rates of obesity, and expanded indications for surgery. It is most commonly performed in individuals over the age of 60, with the peak incidence occurring in the 7th and 8th decades of life, although it is increasingly performed in younger patients due to conditions like avascular necrosis, post-traumatic arthritis, or inflammatory arthropathies. Both men and women undergo THA, with a slight female predominance often attributed to higher rates of osteoarthritis and osteoporosis-related fractures.

Osteoarthritis is the leading indication for THA, accounting for approximately 90\% of all procedures. The perioperative mortality rate for elective THA is very low, typically less than 0.5\%, reflecting advancements in surgical technique, anesthesia, and patient selection. Major morbidity rates, including infection and dislocation, are also low, generally ranging from 1-2\%. The procedure significantly improves quality of life for millions of patients worldwide, making it a public health success story.

\section*{Outcomes \& Prognosis}
The outcomes of total hip replacement are overwhelmingly positive, with the vast majority of patients experiencing substantial relief of pain, significant improvement in hip function, and a marked enhancement in their overall quality of life. Implant survival rates are excellent, with modern prostheses demonstrating 15-year survival rates exceeding 90-95\% and 20-year survival rates often above 80-85\% in national joint registries. The prognosis for appropriately selected patients is generally excellent, with a high likelihood of a durable, pain-free, and functional joint for many years.

Key prognostic factors include patient age (younger patients may have higher revision rates over their lifetime), bone quality, activity level, and adherence to postoperative precautions and rehabilitation. While the replaced joint itself is not susceptible to recurrent arthritis, adjacent joints or the contralateral hip may still develop degenerative changes. The main long-term concerns are aseptic loosening, periprosthetic infection, and polyethylene wear, all of which are uncommon but can necessitate revision surgery. Landmark studies and national joint registries consistently demonstrate the long-term efficacy and cost-effectiveness of THA, solidifying its role as a transformative intervention for end-stage hip disease.

\section*{References}
\begin{enumerate}
  \item American Academy of Orthopaedic Surgeons. Total Hip Replacement. Available at: \url{https://orthoinfo.aaos.org/en/treatment/total-hip-replacement/}. Accessed [Current Year].
  
  \item Schwartz's Principles of Surgery. 11th ed. McGraw-Hill Education; 2019. Chapter 49: Orthopedic Surgery.
  
  \item Sabiston Textbook of Surgery: The Biological Basis of Modern Surgical Practice. 21st ed. Elsevier; 2021. Chapter 64: Orthopedic Surgery.
  
  \item MedlinePlus. Hip Replacement. U.S. National Library of Medicine. Available at: \url{https://medlineplus.gov/hipreplacement.html}. Accessed [Current Year].
\end{enumerate}

\clearpage